\documentclass[11pt]{beamer}
\usetheme{Madrid}
\usecolortheme{default}
\setbeamertemplate{navigation symbols}{}
\setbeamertemplate{footline}{}

\usepackage{amsmath}

\title{Why Does the Messenger Matter More Than the Message?}
\subtitle{Information Channel Effects on Parental Investment in Education}
\author{Jordi Torres Vallverdu}
\institute{MRes TSE --- Behavioral Development Economics}
\date{March 2026}

\begin{document}

\begin{frame}
\titlepage
\end{frame}

%% ---- SLIDE 1: The Puzzle ----
\begin{frame}{The Puzzle: Channel $>$ Content}

\textbf{Fact:} Parents systematically misperceive their children's school performance.
\begin{itemize}
    \item Belief--performance gap $>1$ SD; 1/3 of parents wrong about which child does better (Dizon-Ross, 2019 --- Malawi)
    \item Parents vastly understate missed assignments (Bergman \& Chan, 2021 --- Los Angeles)
\end{itemize}

\vspace{0.3cm}
\textbf{Puzzle:} Providing the \emph{same} information via SMS/WhatsApp instead of report cards generates large effects:
\begin{itemize}
    \item GPA $+0.19$ SD, course failures $-38\%$, at $<$\$63 total cost (Bergman \& Chan, 2021)
    \item Effects concentrate among low-SES households
\end{itemize}

\vspace{0.3cm}
\textbf{Why is this surprising?} Standard theory predicts only the \emph{content} of a signal matters, not its packaging. Report cards are free. Why don't parents use them?

{\small \textit{Neither rational inattention, moral hazard, nor credit constraints fully explain why the delivery channel drives the effect.}}

\end{frame}

%% ---- SLIDE 2: Proposed Explanation & Research Questions ----
\begin{frame}{Proposed Explanation \& Research Agenda}

\textbf{Hypothesis:} The effect operates through \emph{attention and salience}, not (only) belief correction.
\begin{itemize}
    \item Parents under cognitive load do not include ``child performance'' in their active decision model (Gabaix, 2014 --- sparse max)
    \item SMS creates high \emph{contrast} with the prior of ``everything is fine'' $\to$ signal becomes salient (Bordalo, Gennaioli, Shleifer, 2012)
    \item Distinct from wrong beliefs: the variable is not \emph{considered at all}
\end{itemize}

\vspace{0.3cm}
\textbf{How to test it:} Multi-arm RCT that \emph{separates} channel from content:
\begin{itemize}
    \item T1: same info, different channel (SMS) $\to$ identifies attention ($A$)
    \item T2: new info, same channel (paper) $\to$ identifies belief correction ($\lambda$)
    \item T3: new info + SMS $\to$ identifies complementarity
\end{itemize}

\vspace{0.3cm}
\textbf{Research questions:}
\begin{enumerate}
    \item How much of the treatment effect is attention vs.\ belief correction?
    \item Does message salience decay over time (habituation)?
    \item What is the optimal information policy (channel, frequency, targeting)?
\end{enumerate}

\end{frame}

\end{document}
